\section{先选择表示，再选择 Transformer}

上一章解释了为什么架构趋向统一；本章回答统一之前最重要的前置问题：Transformer 到底要建模什么。原始 RGB 视频既连续又冗余，若直接按像素展开，序列长度会压垮训练。Movie Gen 用 Temporal Autoencoder（TAE，时间自编码器）把空间和时间都压缩到 latent（潜变量）空间，再把 latent patch 化为 token。

\subsection{生成模型学习的是哪个随机变量}

“学习视频分布”必须先明确随机变量 $x$ 的定义。它可以是像素、离散 code、连续 latent，或多尺度表示。下面两页从文本对比媒体，说明语言 token 几乎天然适合序列建模，而视频必须先去除大量可预测冗余。

\lecturefigure{slide-177-architecture-questions.jpg}{Movie Gen 架构的三个问题：表示、学习目标与 backbone}{00:11:19}

这三问构成整场课的组织方式：representation 决定 token 长度与信息损失；learning objective 决定网络预测什么；architecture 决定条件如何进入以及计算如何扩展。三者互相制约，不能单独比较“哪个 Transformer block 更强”。

\lecturefigure{slide-178-representation-definition.jpg}{表示学习的起点：若要学习 $P(x)$，必须先定义 $x$}{00:11:43}

文本可通过 tokenizer 映射为离散 token；视频则是 $T_0\times3\times H_0\times W_0$ 的连续张量。若把每个 RGB 值都当 token，数据维度与序列长度远超文本，且相邻像素和帧高度冗余。

\lecturefigure{slide-181-text-vs-media-representation.jpg}{文本表示紧凑离散，媒体表示连续且冗余}{00:13:36}

读图时先比较“语义密度”：一个词 token 可以概括复杂概念，而一个像素只表达局部颜色。媒体表示的目标不是无损压缩全部细节，而是在可接受重建误差下保留生成任务需要的对象、纹理、运动和时间结构。

\begin{importantbox}{表示选择就是计算预算分配}
更强压缩会缩短序列并降低 Transformer 成本，但可能丢失小物体、快速运动和精细纹理；更弱压缩保留信息，却把成本转移到注意力、激活与通信。表示不是数据预处理，而是生成系统的第一层架构。
\end{importantbox}

\subsection{为什么不直接生成像素}

前一节说明媒体冗余，本节把冗余换算为 token 数。设输入视频为
\[
V\in\mathbb R^{T_0\times 3\times H_0\times W_0},
\]
其中 $T_0$ 是帧数，$H_0,W_0$ 是像素高宽，通道数为 3。若直接把每个位置展开，序列规模近似 $T_0H_0W_0$，即使忽略颜色通道也不可行。

\lecturefigure{slide-182-latent-representation.jpg}{直接像素建模概念简单，但高分辨率视频需要复杂上采样或巨量计算}{00:15:52}

Movie Gen 选择连续 latent 而非直接像素。TAE encoder 把输入映射为
\[
X=E(V)\in\mathbb R^{T\times C\times H\times W},
\qquad
\frac{T_0}{T}=\frac{H_0}{H}=\frac{W_0}{W}=8,
\]
其中 $C=16$ 是 latent channels。三个轴各压缩 8 倍，使时空位置数约缩小 $8^3=512$ 倍。

\lecturefigure{slide-185-temporal-autoencoder.jpg}{TAE 将 RGB 视频压缩为连续时空 latent，再由 decoder 恢复像素}{00:19:17}

对于 256 帧、$768\times768$ 视频，压缩后约为 $32\times96\times96$ 个时空位置；若每个位置对应一个 token，则
\[
N=\frac{256}{8}\cdot\frac{768}{8}\cdot\frac{768}{8}
=294{,}912.
\]
论文的 73K context 还包含进一步 patchify，因此 token 数继续下降。这个例子说明“8 倍压缩”必须说明作用在哪些轴上，不能误写成总共只压缩 8 倍。

\begin{knowledgebox}{TAE 如何同时兼容图像与视频}
TAE 在二维空间卷积后增加一维时间卷积，在空间注意力后增加时间注意力。图像被视为单帧视频；通过可变长度编码和丢弃多余解码帧，同一 encoder/decoder 可以处理图像与任意长度视频。
\end{knowledgebox}

\begin{warningbox}{高压缩率不等于免费信息保留}
重建损失低不代表所有生成属性都保留。快速运动、细线、文字、小脸和闪烁纹理可能在 latent 中被平均掉。TAE 必须同时评估空间重建、时间一致性和下游生成质量，不能只报告单帧 PSNR。
\end{warningbox}

\subsection{TAE 的损失与 adversarial loss 边界}

TAE 属于 VAE 路线，但论文还加入 perceptual、discriminator 和 outlier penalty 等目标。对重建 $\hat V=D(E(V))$，可把总体目标写成
\[
\mathcal L_{\mathrm{TAE}}
=\lambda_{\mathrm{rec}}\mathcal L_{\mathrm{rec}}
+\lambda_{\mathrm{perc}}\mathcal L_{\mathrm{perc}}
+\lambda_{\mathrm{adv}}\mathcal L_{\mathrm{adv}}
+\lambda_{\mathrm{KL}}\mathcal L_{\mathrm{KL}}
+\lambda_{\mathrm{OPL}}\mathcal L_{\mathrm{OPL}}.
\]
各 $\lambda$ 是损失权重；重建项约束内容，perceptual 项约束语义视觉相似，adversarial 项鼓励真实感，KL 项正则化 latent，OPL 惩罚异常高范数 latent dots。

\begin{knowledgebox}{第一次出现：adversarial loss}
这里的 adversarial loss 来自一个判别器，用于改善 autoencoder 的解码真实感和压缩能力；它不意味着 Movie Gen 的主生成器是 GAN。主生成器仍以 Flow Matching 训练。Q\&A 中老师强调，加入这类损失可允许 decoder 不必逐像素复制输入，从而支持更高压缩。
\end{knowledgebox}

\teachervoice{01:07:20--01:09:23，老师解释 TAE 的判别器损失来自 VQGAN 路线：像素级 L1 约束太死，adversarial loss 给重建更多感知自由度，并带来额外压缩收益；这与“用 GAN 生成视频”是两回事。}

\begin{lstlisting}[language=Python,caption={视频压缩与 token accounting 的最小核算}]
latent = temporal_autoencoder.encode(video)  # [T/8, C=16, H/8, W/8]
tokens = patchify(latent, patch_t, patch_h, patch_w)
sequence_length = tokens.shape[0]
attention_pairs = sequence_length ** 2
report(sequence_length, attention_pairs, reconstruction_error)
\end{lstlisting}

\subsection{本章小结}

Movie Gen 不直接让 Transformer 建模像素，而是先用 TAE 在时间、高度和宽度上各压缩 8 倍，再 patchify 为 token。压缩将计算变得可行，也引入重建与信息瓶颈。TAE 使用 adversarial loss 不等于主生成器是 GAN；生成 backbone 的学习目标是下一章的 Flow Matching。

